% --- LaTeX Homework Template - S. Venkatraman ---

% --- Set document class and font size ---

\documentclass[letterpaper, 11pt]{article}

% --- Package imports ---

\usepackage{
  amsmath, amsthm, amssymb, mathtools, dsfont,	  % Math typesetting
  graphicx, wrapfig, subfig, float,                  % Figures and graphics formatting
  listings, color, inconsolata, pythonhighlight,     % Code formatting
  fancyhdr, sectsty, hyperref, enumerate, enumitem } % Headers/footers, section fonts, links, lists

% --- Page layout settings ---
\usepackage{bm}
% Set page margins
\usepackage[left=1.35in, right=1.35in, bottom=1in, top=1.1in, headsep=0.2in]{geometry}

% Anchor footnotes to the bottom of the page
\usepackage[bottom]{footmisc}

% Set line spacing
\renewcommand{\baselinestretch}{1.2}

% Set spacing between paragraphs
\setlength{\parskip}{1.5mm}

% Allow multi-line equations to break onto the next page
\allowdisplaybreaks

% Enumerated lists: make numbers flush left, with parentheses around them
\setlist[enumerate]{wide=0pt, leftmargin=21pt, labelwidth=0pt, align=left}
\setenumerate[1]{label={(\arabic*)}}

% --- Page formatting settings ---

% Set link colors for labeled items (blue) and citations (red)
\hypersetup{colorlinks=true, linkcolor=blue, citecolor=red}

% Make reference section title font smaller
\renewcommand{\refname}{\large\bf{References}}

% --- Settings for printing computer code ---

% Define colors for green text (comments), grey text (line numbers),
% and green frame around code
\definecolor{greenText}{rgb}{0.5, 0.7, 0.5}
\definecolor{greyText}{rgb}{0.5, 0.5, 0.5}
\definecolor{codeFrame}{rgb}{0.5, 0.7, 0.5}

% Define code settings
\lstdefinestyle{code} {
  frame=single, rulecolor=\color{codeFrame},            % Include a green frame around the code
  numbers=left,                                         % Include line numbers
  numbersep=8pt,                                        % Add space between line numbers and frame
  numberstyle=\tiny\color{greyText},                    % Line number font size (tiny) and color (grey)
  commentstyle=\color{greenText},                       % Put comments in green text
  basicstyle=\linespread{1.1}\ttfamily\footnotesize,    % Set code line spacing
  keywordstyle=\ttfamily\footnotesize,                  % No special formatting for keywords
  showstringspaces=false,                               % No marks for spaces
  xleftmargin=1.95em,                                   % Align code frame with main text
  framexleftmargin=1.6em,                               % Extend frame left margin to include line numbers
  breaklines=true,                                      % Wrap long lines of code
  postbreak=\mbox{\textcolor{greenText}{$\hookrightarrow$}\space} % Mark wrapped lines with an arrow
}

% Set all code listings to be styled with the above settings
\lstset{style=code}

% --- Math/Statistics commands ---

% Add a reference number to a single line of a multi-line equation
% Usage: "\numberthis\label{labelNameHere}" in an align or gather environment
\newcommand\numberthis{\addtocounter{equation}{1}\tag{\theequation}}

% Shortcut for bold text in math mode, e.g. $\b{X}$
\let\b\mathbf

% Shortcut for bold Greek letters, e.g. $\bg{\beta}$
\let\bg\boldsymbol

% Shortcut for calligraphic script, e.g. %\mc{M}$
\let\mc\mathcal

% \mathscr{(letter here)} is sometimes used to denote vector spaces
\usepackage[mathscr]{euscript}

% Convergence: right arrow with optional text on top
% E.g. $\converge[w]$ for weak convergence
\newcommand{\converge}[1][]{\xrightarrow{#1}}

% Normal distribution: arguments are the mean and variance
% E.g. $\normal{\mu}{\sigma}$
\newcommand{\normal}[2]{\mathcal{N}\left(#1,#2\right)}

% Uniform distribution: arguments are the left and right endpoints
% E.g. $\unif{0}{1}$
\newcommand{\unif}[2]{\text{Uniform}(#1,#2)}

% Independent and identically distributed random variables
% E.g. $ X_1,...,X_n \iid \normal{0}{1}$
\newcommand{\iid}{\stackrel{\smash{\text{iid}}}{\sim}}

% Equality: equals sign with optional text on top
% E.g. $X \equals[d] Y$ for equality in distribution
\newcommand{\equals}[1][]{\stackrel{\smash{#1}}{=}}

% Math mode symbols for common sets and spaces. Example usage: $\R$
\newcommand{\R}{\mathbb{R}}   % Real numbers
\newcommand{\C}{\mathbb{C}}   % Complex numbers
\newcommand{\Q}{\mathbb{Q}}   % Rational numbers
\newcommand{\Z}{\mathbb{Z}}   % Integers
\newcommand{\N}{\mathbb{N}}   % Natural numbers
\newcommand{\F}{\mathcal{F}}  % Calligraphic F for a sigma algebra
\newcommand{\El}{\mathcal{L}} % Calligraphic L, e.g. for L^p spaces

% Math mode symbols for probability
\newcommand{\pr}{\mathbb{P}}    % Probability measure
\newcommand{\E}{\mathbb{E}}     % Expectation, e.g. $\E(X)$
\newcommand{\var}{\text{Var}}   % Variance, e.g. $\var(X)$
\newcommand{\cov}{\text{Cov}}   % Covariance, e.g. $\cov(X,Y)$
\newcommand{\corr}{\text{Corr}} % Correlation, e.g. $\corr(X,Y)$
\newcommand{\B}{\mathcal{B}}    % Borel sigma-algebra

% Other miscellaneous symbols
\newcommand{\tth}{\text{th}}	% Non-italicized 'th', e.g. $n^\tth$
\newcommand{\Oh}{\mathcal{O}}	% Big-O notation, e.g. $\O(n)$
\newcommand{\1}{\mathds{1}}	% Indicator function, e.g. $\1_A$

% Additional commands for math mode
\DeclareMathOperator*{\argmax}{argmax}    % Argmax, e.g. $\argmax_{x\in[0,1]} f(x)$
\DeclareMathOperator*{\argmin}{argmin}    % Argmin, e.g. $\argmin_{x\in[0,1]} f(x)$
\DeclareMathOperator*{\spann}{Span}       % Span, e.g. $\spann\{X_1,...,X_n\}$
\DeclareMathOperator*{\bias}{Bias}        % Bias, e.g. $\bias(\hat\theta)$
\DeclareMathOperator*{\ran}{ran}          % Range of an operator, e.g. $\ran(T) 
\DeclareMathOperator*{\dv}{d\!}           % Non-italicized 'with respect to', e.g. $\int f(x) \dv x$
\DeclareMathOperator*{\diag}{diag}        % Diagonal of a matrix, e.g. $\diag(M)$
\DeclareMathOperator*{\trace}{trace}      % Trace of a matrix, e.g. $\trace(M)$

% Numbered theorem, lemma, etc. settings - e.g., a definition, lemma, and theorem appearing in that 
% order in Section 2 will be numbered Definition 2.1, Lemma 2.2, Theorem 2.3. 
% Example usage: \begin{theorem}[Name of theorem] Theorem statement \end{theorem}
\theoremstyle{definition}
\newtheorem{theorem}{Theorem}[section]
\newtheorem{proposition}[theorem]{Proposition}
\newtheorem{lemma}[theorem]{Lemma}
\newtheorem{corollary}[theorem]{Corollary}
\newtheorem{definition}[theorem]{Definition}
\newtheorem{example}[theorem]{Example}
\newtheorem{remark}[theorem]{Remark}

% Un-numbered theorem, lemma, etc. settings
% Example usage: \begin{lemma*}[Name of lemma] Lemma statement \end{lemma*}
\newtheorem*{theorem*}{Theorem}
\newtheorem*{proposition*}{Proposition}
\newtheorem*{lemma*}{Lemma}
\newtheorem*{corollary*}{Corollary}
\newtheorem*{definition*}{Definition}
\newtheorem*{example*}{Example}
\newtheorem*{remark*}{Remark}
\newtheorem*{claim}{Claim}

\newcommand{\Pp}{\mathbb{P}}
\newcommand{\Dkl}{\mathrm{D}}

% --- Left/right header text (to appear on every page) ---

% Include a line underneath the header, no footer line
\pagestyle{fancy}
\renewcommand{\footrulewidth}{0pt}
\renewcommand{\headrulewidth}{0.4pt}

% Left header text: course name/assignment number
\lhead{Foundations of Machine Learning-- Homework 1}

% Right header text: your name
\rhead{Yixia Yu (yy5091@nyu.edu)}

% --- Document starts here ---

\begin{document}
\section*{A. AdaBoost}

Let $\text{corr}(\mathbf{x}, \mathbf{x}')$ denote the inner product (or unnormalized correlation) of two vectors $\mathbf{x}$ and $\mathbf{x}'$.

\begin{enumerate}
    \item Prove that the distribution vector $(D_{t+1}(1), \ldots, D_{t+1}(m))$ defined by AdaBoost and the vector of components $y_i h_t(x_i)$ are uncorrelated.
    
    \item Fix $\epsilon \in (0, 1/2)$. Let the training sample be defined by $m$ points in the plane with $\frac{m}{4}$ negative points all at coordinate $(1,1)$, another set of $\frac{m}{4}$ negative points all at coordinate $(-1,-1)$, $\frac{m(1-\epsilon)}{4}$ positive points all at coordinate $(1,-1)$, and $\frac{m(1+\epsilon)}{4}$ positive points all at coordinate $(-1,+1)$. Describe the behavior of AdaBoost when run on this sample using boosting stumps, in particular, give the solution the algorithm returns after $T$ rounds.
\end{enumerate}

\section*{B. Boosting algorithms}

\begin{enumerate}
    \item Show that for all $u \in \mathbb{R}$ and integer $p > 1$, $1_{u \leq 0} \leq \Psi_p(-u)$ where $\Psi_p(u) = \max((1+u)^p, 0)$. Show that $\Psi_p$ is convex and differentiable.
    
    \item Use $\Psi_p$ to derive a boosting-type algorithm using coordinate descent. You should give a full description of your algorithm, including the pseudocode, details for the choice of the step and direction, as well as a generalization bound.
\end{enumerate}

\section*{C. Regularized Maxent}

This problem studies $L_2$-regularized Maxent. We will use the notation introduced in class and will denote by $J_S$ the dual objective function minimized given a sample $S$:
\[
J_S(\mathbf{w}) = \frac{\lambda}{2} \|\mathbf{w}\|_2^2 + \mathop{\mathbb{E}}_{x \sim S} \left[ -\log p_{\mathbf{w}}[x] \right],
\]
where $\lambda > 0$ is a regularization parameter. We will assume that the feature vector is bounded: $\|\Phi(x)\|_2 \leq r$ for all $x \in \mathcal{X}$, for some $r > 0$.

\begin{enumerate}
    \item Use McDiarmid's inequality to prove that for any $\delta > 0$, with probability at least $1 - \delta$, the following inequality holds:
    \[
    \left\| \mathop{\mathbb{E}}_{x \sim \mathcal{D}}[\Phi(x)] - \mathop{\mathbb{E}}_{x \sim S}[\Phi(x)] \right\|_2 \leq \sqrt{\frac{2r^2}{m}} \left( 1 + \log \sqrt{\frac{1}{\delta}} \right).
    \]
    
    \item Let $\widehat{\mathbf{w}}$ be the $L_2$-regularized maxent solution for a sample $S$ and $\mathbf{w}_{\mathcal{D}}$ the solution for an infinite sample:
    \[
    \widehat{\mathbf{w}} = \mathop{\text{argmin}}_{\mathbf{w} \in \mathbb{R}^N} J_S(\mathbf{w}) \quad \text{and} \quad \mathbf{w}_{\mathcal{D}} = \mathop{\text{argmin}}_{\mathbf{w} \in \mathbb{R}^N} J_{\mathcal{D}}(\mathbf{w}).
    \]
    where $J_{\mathcal{D}}(\mathbf{w}) = \frac{\lambda}{2} \|\mathbf{w}\|_2^2 + \mathop{\mathbb{E}}_{x \sim \mathcal{D}} \left[ -\log p_{\mathbf{w}}[x] \right]$. Use the definition of $\widehat{\mathbf{w}}$ and $\mathbf{w}_{\mathcal{D}}$ as minimizers (use gradients) to prove that the following inequality holds:
    \[
    \|\widehat{\mathbf{w}} - \mathbf{w}_{\mathcal{D}}\|_2 \leq \frac{\left\| \mathop{\mathbb{E}}_{x \sim S}[\Phi(x)] - \mathop{\mathbb{E}}_{x \sim \mathcal{D}}[\Phi(x)] \right\|_2}{\lambda}.
    \]
    
    \item For any $\mathbf{w}$ and any distribution $\mathcal{Q}$ define $\mathcal{L}_{\mathcal{Q}}(\mathbf{w})$ by $\mathcal{L}_{\mathcal{Q}}(\mathbf{w}) = \mathop{\mathbb{E}}_{x \sim \mathcal{Q}}[-\log p_{\mathbf{w}}[x]]$. Show that
    \[
    \mathcal{L}_{\mathcal{D}}(\widehat{\mathbf{w}}) - \mathcal{L}_{\mathcal{D}}(\mathbf{w}_{\mathcal{D}}) \leq (\widehat{\mathbf{w}} - \mathbf{w}_{\mathcal{D}}) \cdot \left[ \mathop{\mathbb{E}}_{x \sim S}[\Phi(x)] - \mathop{\mathbb{E}}_{x \sim \mathcal{D}}[\Phi(x)] \right] + \frac{\lambda}{2} \|\mathbf{w}_{\mathcal{D}}\|_2^2 - \frac{\lambda}{2} \|\widehat{\mathbf{w}}\|_2^2.
    \]
    
    \item Use that to show that the following inequality holds for any $\mathbf{w}$:
    \[
    \mathcal{L}_{\mathcal{D}}(\widehat{\mathbf{w}}) \leq \frac{1}{\lambda} \left\| \mathop{\mathbb{E}}_{x \sim S}[\Phi(x)] - \mathop{\mathbb{E}}_{x \sim \mathcal{D}}[\Phi(x)] \right\|_2^2 + \mathcal{L}_{\mathcal{D}}(\mathbf{w}) + \frac{\lambda}{2} \|\mathbf{w}\|_2^2.
    \]
    
    \item Conclude by proving that for any $\delta > 0$, with probability at least $1 - \delta$, the following inequality holds:
    \[
    \mathcal{L}_{\mathcal{D}}(\widehat{\mathbf{w}}) \leq \inf_{\mathbf{w} \in \mathbb{R}^N} \mathcal{L}_{\mathcal{D}}(\mathbf{w}) + \frac{\lambda}{2} \|\mathbf{w}\|_2^2 + \frac{2r^2}{\lambda m} \left( 1 + \sqrt{\log \frac{1}{\delta}} \right)^2.
    \]
\end{enumerate}

\section*{D. Randomized Halving}

In class, we showed that, in the realizable scenario (at least one expert is always correct), the number of mistakes made by Halving is upper bounded by $\log_2 N$. Here, we consider for the same realizable scenario a randomized version of Halving defined as follows.

As for Halving, let $H_t$ denote the set of remaining experts at the beginning of round $t$, with $H_1 = H$ the full set of $N$ experts. At each round, let $r_t$ be the fraction of experts in $H_t$ predicting 1. Then, the prediction $\widehat{y}_t$ made by the algorithm is 1 with probability
\[
p_t = \left[ \frac{1}{2} \log_2 \frac{1}{1-r_t} \right] 1_{r_t \leq \frac{3}{4}} + 1_{r_t > \frac{3}{4}},
\]
0 with probability $1 - p_t$. The true label $y_t$ is then received and $H_{t+1}$ is derived from $H_t$ by removing all experts who made a mistake.

\begin{enumerate}
    \item Write the pseudocode of the algorithm.
    
    \item Define the potential function $\Psi_t = \log_2 |H_t|$. Let $\mu_t = 1_{y_t \neq \widehat{y}_t}$, prove that for all $t \geq 1$, $\mathbb{E}[\mu_t] \leq \frac{\Psi_t - \Psi_{t+1}}{2}$.
    
    \item Show that the expected number of mistakes made by randomized Halving is at most $\frac{1}{2} \log_2 N$.
    
    \item (Bonus question) Prove that no randomized algorithm makes fewer than $\lfloor \frac{1}{2} \log_2 N \rfloor$ mistakes, in expectation.
\end{enumerate}

\end{document}